This report documents the development and implementation of an investment
recommendation system for FinPUC, a fictitious financial advisory platform
that generates personalized portfolios for clients in the US stock market.
The system serves five risk profiles and pursues a dual objective: maximizing
client expected returns and maximizing firm revenue from transaction fees.

The investment universe comprises 601 stocks filtered from 1,406 tickers.
Three methodological layers of increasing complexity were implemented and
compared: Markowitz mean-variance optimization with semiannual rebalancing,
Black-Litterman with four quantitative view variants (general momentum,
6-month and 1-year capitalization-weighted momentum, and an assumed
macroeconomic unemployment scenario), and a Monte Carlo simulation of 5,000
trajectories over 260 weeks modeling client acceptance, withdrawal, and firm
commissions across six techniques.

Results show that the Black-Litterman variant with 6-month
capitalization-weighted momentum dominates the recommender ranking, achieving
an expected terminal wealth of USD~8,974 and a P4 score of 9,417 for the
most aggressive risk profile with pandemic training data. Including
pandemic data (2020 - 2022) in training deteriorates aggressive Markowitz
portfolios by up to 1.1 Sharpe points, while Black-Litterman variants
exhibit greater stability across regimes.

\textbf{Keywords:} portfolio optimization, Markowitz, Black-Litterman, Monte
Carlo simulation, recommender system, semiannual rebalancing, CVXPY.